\section{Methods}
\label{sec:methods}

\begin{figure*}[t]
\centering
\includegraphics[width=\textwidth]{figures/fig1_workflow.pdf}
\caption{Study design. For every held-out participant $p$, all preprocessing statistics, model fitting and channel ranking use only the other participants $R_{-p}$. Rankings are cross-fitted: $R_{-p}$ is split into five participant-disjoint inner folds, and attributions are computed only on inner-held-out participants. Subsets of $k$ channels are then retrained from scratch and tested on $p$. Decision rules (margin, statistical unit, bootstrap, multiplicity, $k^*$) were frozen before any sparse-channel result existed.}
\label{fig:workflow}
\end{figure*}

\subsection{Datasets}

\paragraph{GRABMyo (primary)}
GRABMyo v1.1.0~\cite{pradhan2022multiday} contains recordings from \GRABMyoParticipants\ able-bodied participants (\GRABMyoMale\ male, \GRABMyoFemale\ female; age \GRABMyoAgeMean\ $\pm$ \GRABMyoAgeSD\ years) on \GRABMyoSessions\ different days. Signals were recorded with an EMGUSB2+ amplifier (OT Bioelettronica) at 2048~Hz. In each session participants performed 16 hand and wrist gestures plus rest, with 7 trials of 5~s each. We analysed the 16 forearm electrodes. These form two circumferential rings of 8 electrodes (F1--F8 and F9--F16), 2~cm apart along the forearm, with electrode $j$ of one ring aligned with electrode $j+8$ of the other. Mean circumference at the ring position was \GRABMyoCircMean\ $\pm$ \GRABMyoCircSD~cm, giving about \GRABMyoArcSpacing~cm between neighbouring electrodes in a ring. The recordings are monopolar, so the selection unit in the primary analysis is a single electrode. All \GRABMyoFilesVerified\ distributed files matched the dataset's published SHA-256 checksums. All \GRABMyoRecords\ records were processed, with \GRABMyoRecordsFailed\ failures. The primary task is the 16 active gestures (rest excluded), for which chance Macro-F1 is about 0.063.

\paragraph{NinaPro DB5 (exploratory methodological replication)}
DB5~\cite{pizzolato2017comparison} provides \DBFiveSubjects\ able-bodied subjects (\DBFiveMale\ male, \DBFiveFemale\ female; age \DBFiveAgeMean\ $\pm$ \DBFiveAgeSD\ years, from the files' metadata) recorded with two Myo armbands (16 channels, 200~Hz). We used exercise~1 (12 finger movements, 6 repetitions each), segmented with the refined \texttt{restimulus}/\texttt{rerepetition} labels. The DB5 values are 8-bit integer codes, not physical units. The spectra show no powerline peak, and their power extends to the Nyquist frequency. DB5 was analysed separately and never pooled with GRABMyo. Its channel numbers bear no anatomical correspondence to GRABMyo's, so only the behaviour of the procedure is compared across datasets.

\subsection{Preprocessing and features}
GRABMyo signals were band-pass filtered (20--450~Hz, 4th-order zero-phase Butterworth) and notch filtered at 50~Hz ($Q = 30$). A spectral audit chose 50~Hz over 60~Hz: mean band-power ratio was 10.35 at 50~Hz and 0.62 at 60~Hz. Windows were 250~ms with 50\% overlap and never crossed trial boundaries. From each channel we extracted ten features: mean absolute value (MAV), root mean square, variance, waveform length, zero crossings (ZC), slope sign changes (SSC), integrated EMG (IEMG), Willison amplitude (WAMP), mean frequency and median frequency. This gives a 160-dimensional vector per window, \GRABMyoWindows\ windows in total and \GRABMyoActiveWindows\ from active gestures. The ZC, SSC and WAMP amplitude threshold was set to $10^{-2}$~mV, about 1.15 times the median absolute first difference of the signal. This single label-free scalar was computed from the pooled recordings, not per fold. That is a minor departure from strict fold-wise estimation, which we note as a limitation. An earlier threshold 87 times smaller had saturated WAMP and was discarded. For DB5 we used a 10--95~Hz band-pass, no notch, 250~ms windows (50 samples) and a threshold of 7 ADC units, derived in the same way.

\subsection{Classifiers and leave-one-subject-out evaluation}
All evaluation is leave-one-subject-out (LOSO): 43 folds for GRABMyo and 10 for DB5. In fold $p$, z-scoring statistics, model fitting and channel ranking use only $R_{-p}$. The test participant contributes only to the score. XGBoost used 300 trees, depth 6, learning rate 0.1, row and column subsampling of 0.8, and the \texttt{hist} method. It was trained on all \FoldTrainingWindows\ outer-training windows. For compute reasons, the RBF support vector machine ($C = 1$, $\gamma$ = \texttt{scale}) used a stratified 40,000-window subsample and the random forest (200 trees, minimum leaf size 2) a 200,000-window subsample. On one calibration fold, halving the SVM subsample from 40,000 to 20,000 windows lowered Macro-F1 from \SVMCalibFortyK\ to \SVMCalibTwentyK. The SVM had therefore not saturated, and its result is a lower bound. No hyperparameter was tuned on any outer test participant. The primary metric is participant-level Macro-F1 over the 16 active gestures; accuracy and balanced accuracy were also recorded. XGBoost was prespecified as the primary model for the channel analyses on grounds of compute.

\subsection{Cross-fitted channel attribution and the SR-CFA score}
Within each outer fold, $R_{-p}$ was split into $K = 5$ participant-disjoint inner folds. For inner fold $j$, an XGBoost model was trained on the other inner folds (capped at a 200,000-window stratified subsample and z-scored with inner-training statistics only). Exact TreeSHAP attributions~\cite{lundberg2020local} were then computed for up to 20,000 windows of the inner-held-out participants. Per-window channel importance averaged absolute attributions over the 16 classes and the 10 features of the channel,
\begin{equation}
\Phi_{i,c} = \frac{1}{G F}\sum_{g=1}^{G}\sum_{f=1}^{F} \lvert \phi_{i,g,c,f} \rvert ,
\end{equation}
and was averaged over windows and normalised to sum to one across channels, giving $I_{j,c}$. Across the five inner folds, $A_c = \operatorname{median}_j I_{j,c}$ measures attribution magnitude. $U_c = \operatorname{median}_j \lvert I_{j,c} - A_c \rvert$ measures its instability across disjoint groups of training participants. Channels were ranked by
\begin{equation}
Q_c(\lambda) = A_c - \lambda\, U_c , \qquad \lambda \in \{0, 0.25, 0.5, 1\}.
\label{eq:srcfa}
\end{equation}
$\lambda = 0$ is cross-fitted SHAP magnitude; $\lambda > 0$ is the stability-regularised score (SR-CFA). A leakage assertion verified in every fold that no inner-held-out participant appeared in its inner training set. Unit tests verified that a participant-specific signal planted in the data could not raise that participant's score.

The frozen plan required $\lambda$ to be chosen on outer-training data only but did not specify a criterion, and no selection procedure was implemented. We therefore select no $\lambda$ and report all four (a post-hoc specification, Section~\ref{sec:deviations}).

\subsection{Sparse sensor curve and non-inferiority (H1)}
For every outer fold, $\lambda$ and $k \in \{1,\ldots,16\}$, the top-$k$ channels of the fold's own ranking were used to retrain XGBoost from scratch. Each model used the fold's fixed 200,000-window stratified outer-training subsample and was tested on the held-out participant. The $k = 16$ model of the same curve is the reference, so every paired difference $\Delta_{p,k} = \mathrm{F1}_{p,k} - \mathrm{F1}_{p,16}$ compares models with identical training budgets. Non-inferiority at margin $\delta$ requires the lower bound of the 95\% participant-cluster bootstrap interval of $\bar\Delta_k$ to exceed $-\delta$. We define $k^* = \min\{k < 16 : \mathrm{LCB}_{95}(\bar\Delta_k) > -\delta\}$, recorded as ``none'' if no $k$ qualifies. The primary margin $\delta = 0.02$ Macro-F1 is an engineering tolerance, not a clinically validated threshold. $\delta = 0.01$ and $0.03$ are prespecified sensitivity margins.

\subsection{Comparison with other selection methods (H2)}
At $k \in \{4, 6, 8, 12\}$ we compared six rankings, each computed on outer-training data only:
\begin{itemize}
\item SR-CFA ($\lambda = 0.5$);
\item cross-fitted SHAP ($\lambda = 0$);
\item permutation importance, from a model fitted on a 20,000-window subsample and evaluated on 8,000 further windows with 3 repeats;
\item mutual information, estimated on the same subsample;
\item a geometric baseline, which selects electrodes by farthest-point sampling on the $2 \times 8$ montage;
\item random subsets: 100 draws per participant at each $k$, and 500 at the prespecified $k^*$.
\end{itemize}
The value $\lambda = 0.5$ for this comparison had no recorded rationale. Because the P4 comparison includes $\lambda = 0$ as well, both rankings are reported. Drawing hundreds of random subsets per participant was only tractable with a lighter comparison model (XGBoost with 50 trees, trained on a 150,000-window stratified subsample), applied identically to every method. Its absolute scores are therefore lower than those of the sparse curve and are compared only within the comparison. H2 contrasts each attribution ranking with the participant's mean over its random subsets (Holm over 8 tests). We also report where the SR-CFA subset falls within each participant's random distribution.

\subsection{Stability regularisation (H3) and faithfulness (H4)}
H3 compared each $\lambda > 0$ with $\lambda = 0$ on the sparse curve at $k \in \{4, 6, 8, 12\}$ (Holm over 12 tests). For H4 we used deletion with retraining~\cite{hooker2019benchmark}. For $n \in \{2, 4, \ldots, 14\}$ retained channels, the model was retrained after removing either the $16 - n$ highest-ranked channels (keeping the $n$ lowest) or the $16 - n$ lowest-ranked (keeping the top $n$). Random removal is the mean over the random subsets above, which exist for $n \in \{4, 6, 8, 12\}$. All three arms used the same comparison model and training subsample. As a check on protocol equivalence, the top-$n$ arm selects the same electrodes as the H2 cross-fitted SHAP subsets in all \FaithProtocolCells\ matching participant-by-$k$ cells. Its Macro-F1 agrees with them to within a mean absolute difference of \FaithProtocolMeanDiff\ (maximum \FaithProtocolMaxDiff). The residual differences are numerical: z-scoring was applied per subset rather than once over all features. A faithful ranking gives top-removed $<$ random $<$ bottom-removed. Paired contrasts were Holm-corrected within H4.

\subsection{Statistics}
The participant is the statistical unit ($n = 43$; $n = 10$ for DB5). Windows are never treated as independent. Confidence intervals are 95\% percentile intervals from 10,000 cluster-bootstrap resamples of participants. Paired comparisons use the Wilcoxon signed-rank test with the matched-pairs rank-biserial correlation $r_{rb}$ as effect size. Holm correction is applied within each confirmatory family. Exploratory analyses are uncorrected and labelled as such. No participant was excluded.

\subsection{Exploratory analyses}
\paragraph{Ranking stability} We measured the agreement of rankings computed from single participants' attributions, from the five inner folds, and from the 43 outer folds. Outer folds share about 98\% of their training participants, so their agreement is only an upper bound.

\paragraph{Cross-session} XGBoost (full configuration) was trained on session~1 of $R_{-p}$ (CS1) and tested on each session of $p$. It was also trained on sessions 1 and 2 (CS2) and tested on session 3 of $p$. CS2 has roughly twice the training data of CS1, so the two conditions differ in both session diversity and quantity.

\paragraph{Bipolar montage} To express the selection unit as a conventional bipolar sensor, we derived $B_j = F_j - F_{j+8}$ ($j = 1, \ldots, 8$) and re-extracted features. The count-feature threshold was re-derived for this signal ($5 \times 10^{-3}$~mV), because the bipolar signal is about half the monopolar amplitude. We then repeated the LOSO baseline, the cross-fitted ranking and the sparse curve, with the $k = 8$ model as reference.

\paragraph{DB5} The LOSO baseline, cross-fitted ranking ($K = 5$) and sparse curve were repeated on DB5, using all outer-training windows and the $k = 16$ model as reference.

\paragraph{Fixed layout} In a post-hoc analysis, a single 12-electrode layout was derived from a cross-fitted ranking ($\lambda = 0.5$) on the 22 odd-numbered participants. Models with 16 and with 12 electrodes were trained on those participants and tested on the 21 even-numbered participants. This is one split, not a LOSO estimate.

\paragraph{Feature redundancy} IEMG equals 512 times MAV for these windows. We repeated the comparison-model LOSO with IEMG removed and compared both Macro-F1 and channel attribution.

\subsection{Deviations from the analysis plan}
\label{sec:deviations}
The following changes were made after results existed. None alters a prespecified decision rule.
\begin{enumerate}
\item $\lambda$ was not selected; see above.
\item The H2--H4 families were specified in detail (tests and $k$ sets) after the sparse curve and selection-baseline results had been computed.
\item The comparison model for H2 and H4 was reduced from 150 to 50 trees for compute reasons before any baseline result existed.
\item An earlier faithfulness run mixed arms trained with different budgets and was discarded, and all arms were recomputed under one protocol.
\item An earlier bipolar run used a count threshold inherited from the monopolar signal, and was discarded and recomputed.
\item The fixed-layout analysis was not planned.
\item The window-length sensitivity and the 17-class task, both listed as exploratory in the plan, were not run.
\end{enumerate}
